\SourceLeaf{161}$m, m', m'', m'''$ etc. selbst, so doch ihr Verhältnis kennen (d. h. das Verhältnis der Zahlen, durch die sie ausgedrückt werden, wobei bei heterogenen Größen für jede eine beliebige Einheit zugrunde gelegt wird). Mit gesundem Menschenverstand wird man fast immer wenigstens eine Möglichkeit finden, sie einigermaßen abzuschätzen. Wo aber jede Kenntnis gänzlich fehlt, liegt es wahrhaftig nicht an der Methode selbst, dass sie ihre Dienste versagen muss. Auf \textls{Sachkenntnis} kann bei der Anwendung der Methode der kleinsten Quadrate niemals verzichtet werden. Es wird also jedes Mal von der \textls{Sachkenntnis} abhängen, inwieweit man bei \textls{Längenmessungen} den zu erwartenden mittleren Fehler für kleinere und größere Entfernungen gleich oder ungleich ansetzen soll und im letzteren Fall, wie man ihn von der Größe abhängig machen soll. Betrachten wir zum Beispiel kleine Entfernungen. Misst man etwa mit einer 3 Meter langen Messlatte, so hängt der Fehler, wenn die Entfernung selbst nur \textls{sehr wenige} Meter beträgt, ganz oder fast allein von der Teilung ab. Muss aber der Maßstab viele hundert Mal umgeschlagen werden, so kommt eine ganz andere Fehlerquelle ins Spiel.

\begin{center}
. . . . . . . . . . . .
\end{center}
